E1 hoàn thành 600/600 lượt trên 15 cấu hình. Hình~\ref{fig:e1rt} cho thấy runtime tăng cùng kích thước trong ba họ đồ thị. Hình~\ref{fig:e1work} đối chiếu số lượt xử lý đỉnh. Bảng~\ref{tab:e1} tóm tắt các cấu hình 100.000 đỉnh. LPA và FLPA có cùng bước tính phiếu. Bên cạnh số lượt, phân bố bậc tại các đỉnh được ghé và lần quét để đưa lại láng giềng cũng ảnh hưởng thời gian.

Trên ER 100.000 đỉnh, LPA dùng trung bình 44,30 lượt trên một đỉnh, còn FLPA dùng 29,84. Thời gian trung bình lần lượt là 10,190 và 6,435 giây. Trung bình tỷ lệ ghép cặp LPA/FLPA là $1.58\pm0.69$. Ở mỗi seed, cả hai trả về một cụm lớn cùng một nhãn riêng cho từng đỉnh cô lập. Vì vậy số cụm trung bình 6,3 chủ yếu phản ánh chính sách giữ isolate. Retention-LPA dừng ở trung bình 0,831 giây và trả hơn 31.000 nhóm. Tốc độ của retention đi cùng một nghiệm rất mịn, cũng là xu hướng tác giả quan sát trên ER.

Trên BA 100.000 đỉnh, LPA và FLPA cùng trả một cụm ở cả 10 lần, với thời gian lần lượt $6.679\pm2.780$ và $5.223\pm1.672$ giây. Trung bình tỷ lệ ghép cặp là $1.30\pm0.34$. FLPA giảm lượt đỉnh từ 15,73 xuống 9,62 lượt trên đỉnh. Quét lại láng giềng khi cập nhật hub vẫn có chi phí đáng kể. Retention-LPA tạo khoảng 27.403 cụm trung bình. Leiden thu được trung bình 21,9 cụm và $Q=0.295$. Do BA thiếu nhãn chuẩn, số cụm và modularity cho thấy các dạng phân hoạch khác nhau và làm rõ ảnh hưởng của mục tiêu tối ưu.

SBM 100.000 đỉnh có bậc trung bình thực từ 9,984 đến 10,025 và phần cạnh ngoại khối từ 0,5994 đến 0,6013. Tỷ lệ thời gian LPA/FLPA là $1.40\pm0.54$ theo cặp. LPA và FLPA lần lượt có khoảng 43,18 và 32,83 lượt xử lý trên mỗi đỉnh. Cả hai có NMI xấp xỉ 0 và ACC xấp xỉ 0,01 tại $\mu=0.60$. Retention-LPA cho NMI 0,508 trong khi ARI là 0,0015 và ACC là 0,0059 với hơn 31.000 cụm. Các mảnh nhỏ mang thông tin về nhóm gốc theo NMI, nhưng rất ít cặp cùng nhóm được giữ và các mảnh dư bị phạt khi ghép một-một. Leiden đạt modularity 0,303 trên cấu hình này và NMI 0,021. Hai mục tiêu chất lượng đánh giá những đặc tính khác nhau của phân hoạch.

Bảng~\ref{tab:paper} đưa ra thời gian Figure~1 tại 50.000 đỉnh và độ lệch tương đối của bản Python. Gọi $\bar t_P$ là mean đo được bằng Python và $\bar t_R$ là mean tham chiếu từ hình hoặc phép đo C. Cột sai lệch dùng phần trăm
\begin{equation}
  \Delta = 100\left(\frac{\bar t_P}{\bar t_R}-1\right).
  \label{eq:deviation}
\end{equation}
Điểm tham chiếu từ hình được làm tròn đến 0,05 giây. Độ lệch tương đối dễ nhạy với sai số đọc khi $\bar t_R$ nhỏ. Panel tăng kích thước cắt một số điểm lớn, nên so sánh SBM 100.000 đỉnh tại $\mu=0.60$ sử dụng panel thay đổi $\mu$. Ở panel đó, thời gian LPA và FLPA xấp xỉ 5,80 và 2,70 giây, tỷ số khoảng 2,15. Thử nghiệm Python cho tỷ số của hai mean là 1,44 và mean các tỷ số ghép là 1,40. Cả ba họ đều cho FLPA nhanh hơn LPA theo trung bình ghép cặp ở tất cả kích thước, trong khoảng 1,20–2,46 lần. Độ lớn cải thiện khác nhau theo họ và kích thước. Mốc trên 700 lần trong bài báo thuộc mạng us-patents, được báo cáo ở Bảng~2~\cite{flpa}.

Bảng~\ref{tab:author} bổ sung phép đo trên cùng topology bằng mã C. Cả 54 lượt C hoàn thành ở hai kích thước 10.000 và 100.000. Tại 100.000 đỉnh, tỷ số mean LPA/FLPA của C là 1,98 trên ER, 2,50 trên BA và 2,39 trên SBM. Trên đúng ba cặp seed, thời gian FLPA Python/C có tỷ số mean lần lượt 4,49, 12,93 và 18,06. Ở SBM, FLPA của cả hai backend đều có NMI xấp xỉ 0,0001. Retention-LPA ở cả hai backend trả khoảng 31.000 cụm và có NMI gần 0,509. Hai kết quả này cho thấy hiện tượng gộp cụm và phân mảnh cũng xuất hiện ở mã tham chiếu cùng đồ thị.

Phép so sánh C giữ nguyên cạnh và seed số trên cùng máy Apple M1. Hai backend có thể khác RNG, thứ tự duyệt, cấu trúc lưu trữ và chi phí bộ đếm. Chênh lệch với Figure~1 còn chịu ảnh hưởng của Xeon E5-2697, số lần lặp 10 và 2.000, generator và phiên bản thư viện. Khi đánh giá độ lớn của lợi ích FLPA, cần đọc cả độ lệch chuẩn, số lượt xử lý và kiểu phân hoạch mà mỗi thuật toán tạo ra.
